\scanpage{39}
i.e.
\[
\mu\subset(\mathrm{dét}\,M)^*\qquad\mu=-\mu,\quad\mathrm{card}\,\mu=2,
\]
le \deleted{\unsure{l'ens}} polyèdre régulier $SF_{M,\mu}$ admet un iso. can.
\[
\underline{\mathrm{Or}}(SF_{M,\mu})\simeq\mu\tag*{(91)}
\]
Pour $\Lambda=\mathbb Z$, $(\mathrm{det}\,M)^*$ est déjà réduit à deux éléments, et \deleted{\unsure{$P$}} $P_M$ est un 2-polyèdre régulier connexe, dont l'ens des orientations est en bijection can. avec $\mu=(\mathrm{det}\,M)^*$. Pour $\Lambda=\mathbb Z/n\mathbb Z$, avec $n\geqslant3$, alors $(\mathrm{det}\,M)^*/\{\pm1\}$ est de cardinal $\varphi(n)/2$, donc $P_M$ a $\varphi(n)/2$ composantes connexes. Il est connexe sss $\varphi(n)\ne2=1$, i.e.
\[
P_M\ \text{connexe}\iff\varphi(n)=2\iff n\in\{3,4,6\},\tag*{(92)}
\]
le premier cas où $P_M$ est disconnexe étant relatif à la valeur $n=5$, où on trouve pour $P_M$ un ``biicosaèdre''. Il faudrait revenir sur ce cas par la suite (ainsi que \deleted{sur tous} \inserted{sur chacun} des cas $n\in\{3,4,5,6\}$). Pour le moment, je voudrais, dans le cas $\Lambda$ fini i.e. $\Lambda\simeq\mathbb Z/n\mathbb Z$, et dans le cas $P_M$ disconnexe i.e. $n=5$ ou $n\geqslant7$, dégager des éléments de structure supplé%
% archive sheet number: 38
